\section{Conclusion}

Splitting the per-packet cost of an encrypted data plane into a fixed term and a
per-byte term, and reporting their ratio, gives a quantity that survives the
move from one machine to another where neither term alone does. On the host
measured here, \bestCipher\ gives $a = \bestA$ cycles, $b = \bestB$ cycles per
payload byte, and $s^{*} = \bestSStar$ bytes against a maximum carriable payload
of \protoMaxPayload\ bytes.

The same decomposition was the useful lens at the further boundaries this study
could reach. At the kernel boundary, a submission path that issues no system
calls measures directly what batching curves could previously only bound, and
this run shows how sharply that measurement depends on the poller having a core
of its own. At the isolation boundary, the tiers that built cost the same number
of boundary crossings and differed by less than the run-to-run spread, which is
a result about the host rather than about isolation. At the bus, the fixed cost
and the transfer term together decide a batch size below which an accelerator
cannot win regardless of how fast its kernel is.

Two methodological points carried more weight than expected. Measuring two
independent implementations of each cipher separated what the algorithm costs
from what one library's version costs, and it made the masked arm meaningful
rather than a comparison between a software path and a hardware one. And writing
every control's prediction down before running it turned the controls into
something that could fail. One of them did, at first, because its threshold
tested the wrong quantity, and fixing the control rather than the threshold is
what made the instrumentation claim honest.
